\Needspace{16\baselineskip}
\section{附加实验：性能提升与问题改进}
\subsection{增强、正则化、加宽与集成}
B1以CV最优配置为起点，固定45000/5000划分、种子0和40轮预算，逐项加入增强、平滑与加宽，最后替换架构，结果见表\ref{tab:b1}与图\ref{fig:b1}。
\begin{table}[H]\centering\small
\caption{B1累计配方及架构替换，种子0、40轮。集成成员按验证表现选择。}\label{tab:b1}
\ReportTable{tables/b1_ablation}
\end{table}
\begin{figure}[H]\centering
\includegraphics[width=\linewidth]{b1_ablation.pdf}
\caption{B1相邻配方的验证准确率变化（种子0）。虚线分隔架构替换；绝对准确率与参数量见表\ref{tab:b1}。}\label{fig:b1}
\end{figure}

裁剪翻转将测试准确率从88.44\%提高至92.64\%，收益最大；它扩展了训练图像的空间变化。Cutout和标签平滑进一步改变正则化方式，准确率小幅波动。加宽VGG后，参数量由约2.20M增至6.68M，测试准确率达到93.87\%。ResNet-56约0.86M参数，继承VGG选出的优化配方，测试为91.63\%，低于加宽VGG。

按末轮验证准确率选出的三个成员为加宽VGG、标签平滑VGG和裁剪翻转VGG。三模型概率平均达到94.36\%测试准确率，比加宽单模型高0.49个百分点，集成综合了不同配方的预测。

\FloatBarrier
\Needspace{7\baselineskip}
\subsection{深度退化与残差连接}
对Plain和ResNet使用相同主路径深度、BN、裁剪翻转、峰值学习率0.1，训练30轮。图\ref{fig:depth}和表\ref{tab:depth}显示，Plain从20层增至56层时，训练子集误差由5.38\%升至17.20\%，验证误差由10.76\%升至19.58\%。训练与验证误差同时上升，呈现深度增加后的优化退化。
\begin{figure}[H]\centering
\includegraphics[width=\linewidth]{depth_train_error.pdf}
\caption{Plain与ResNet的训练误差曲线（种子0）。}\label{fig:depth}
\end{figure}
\begin{table}[H]\centering\small
\caption{不同深度的末轮错误率（\%），种子0。}\label{tab:depth}
\begin{tabular}{lcccc}
\toprule
层数 & Plain 训练误差 & Plain 验证误差 & ResNet 训练误差 & ResNet 验证误差 \\
\midrule
20 & 5.38 & 10.76 & 3.62 & 8.78 \\
32 & 7.68 & 11.76 & 2.80 & 7.76 \\
56 & 17.20 & 19.58 & 2.46 & 7.62 \\
\bottomrule
\end{tabular}

\end{table}

ResNet的相应训练误差由3.62\%降至2.46\%，验证误差由8.78\%降至7.62\%，支持残差捷径缓解深度退化\cite{he2016resnet}。残差捷径提供较直接的梯度传播路径，使深层网络更容易优化。

\FloatBarrier
\Needspace{8\baselineskip}
\subsection{置信度失真与温度缩放}
表\ref{tab:calibration}同时保留正确率、NLL和ECE。无增强基础模型平均置信度94.65\%，而准确率88.44\%，显示过度自信。Cutout原始ECE为1.98\%，验证集拟合$T=1.26$后测试ECE降至0.54\%，NLL从0.214降至0.205，而准确率仍为93.01\%。图\ref{fig:reliability}提供分箱观察。
\begin{table}[H]\centering\small
\renewcommand{\arraystretch}{1.0}
\caption{测试集校准结果（种子0）。TS为温度缩放，Aug.为裁剪翻转，LS为标签平滑；Acc./Conf./ECE单位为\%。}\label{tab:calibration}
\begin{tabular}{lccccccc}
\toprule
模型 & Acc. & Conf. & NLL & ECE & $T$ & NLL(TS) & ECE(TS) \\
\midrule
Base & 88.44 & 94.65 & 0.441 & 6.26 & 1.69 & 0.356 & 1.69 \\
Aug. & 92.64 & 96.09 & 0.251 & 3.47 & 1.55 & 0.215 & 0.84 \\
Cutout & 93.01 & 94.92 & 0.214 & 1.98 & 1.26 & 0.205 & 0.54 \\
LS & 92.84 & 84.97 & 0.281 & 7.87 & 0.63 & 0.222 & 1.52 \\
Wide & 93.87 & 85.90 & 0.254 & 8.01 & 0.63 & 0.192 & 1.25 \\
ResNet56 & 91.63 & 85.43 & 0.312 & 6.20 & 0.71 & 0.271 & 2.01 \\
\bottomrule
\end{tabular}

\end{table}
\begin{figure}[H]\centering
\includegraphics[width=.92\linewidth]{reliability.pdf}
\caption{可靠性图（15等宽箱）。上排为箱内平均置信度与准确率，下排为样本占比；虚线表示理想校准。}\label{fig:reliability}
\end{figure}

\textbf{标签平滑模型表现为欠自信。}LS模型置信度84.97\%低于准确率92.84\%（ECE为7.87\%）；加宽平滑模型也欠自信。平滑目标降低了正确类别的目标概率。两者拟合$T\approx0.63$后，logits差异被放大，ECE分别降至1.52\%、1.25\%。

\FloatBarrier
\Needspace{7\baselineskip}
\subsection{平移敏感性与抗混叠池化}
测试对$(d_y,d_x)\in[-4,4]^2$的81种整数位移进行评估，采用标准化零填充和边界截断。MaxBlurPool先做步幅1的2$\times$2最大池化，再用$[1,2,1]^\mathsf{T}[1,2,1]/16$的3$\times$3逐通道核低通滤波、反射padding并以步幅2下采样。

设图像$i$在81个位置中被预测为类别$c$的次数为$n_{i,c}$，一致性为不同位置对预测相同的比例，再对图像取平均：
\begin{equation}
 C=\frac1n\sum_{i=1}^n\frac{\sum_c n_{i,c}(n_{i,c}-1)}{81\times80}.
\end{equation}

\begin{table}[H]\centering\small
\caption{池化与增强的平移评估（种子0，40轮）。平均包含零位移，共81个位置。}\label{tab:shift}
\ReportTable{tables/shift}
\end{table}
\begin{figure}[H]\centering
\includegraphics[width=\linewidth]{shift_accuracy.pdf}
\caption{按$\max(|d_y|,|d_x|)$分组的平均准确率。半径0至4分别含1、8、16、24、32个位移。}\label{fig:shift}
\end{figure}

\textbf{抗混叠主要改善扰动稳定性。}无增强时平移平均准确率由83.06\%升至85.88\%，与裁剪翻转结合后也保持收益（表\ref{tab:shift}）。低通滤波减轻了下采样对小幅平移的敏感性。
